\chapter{Таблица команд}\label{app:komandy}

Все команды дрона. Дрон \kod{d} получен так: \kod{d = show.drone(N)}. Статусы: \status{проверено} — выполнялось на живом дроне и загружается; \status{симулятор} — работает на экране, на живом дроне не проверено.

\begin{center}
\small
\renewcommand{\arraystretch}{1.3}
\begin{longtable}{@{}p{4.9cm}p{6.6cm}p{2.5cm}@{}}
\toprule
Запись & Что делает & Статус \\
\midrule
\endhead
\kod{calibrate(seconds)} & калибровка перед взлётом, обычно первый шаг & проверено \\
\kod{takeoff(z, seconds)} & взлёт на высоту \kod{z} см за \kod{seconds} с & проверено \\
\kod{rgb(color)} & цвет светодиода: имя, номер 0–7 или \kod{(r, g, b)}; занимает 0\,с & проверено \\
\kod{rgb_value(r, g, b)} & цвет по трём каналам 0–255; то же, что \kod{rgb((r, g, b))} & проверено \\
\kod{code(cell, seconds)} & лететь к клетке по прямой за \kod{seconds} с & проверено \\
\kod{code_s(cells, seconds)} & пролететь клетки списка или строки \kod{'12 23 27'}, время делится поровну & проверено \\
\kod{set_z(z, seconds)} & перейти на высоту \kod{z}; на той же высоте — ожидание & проверено \\
\kod{code_z(cell, z, seconds)} & лететь к клетке и на итоговую высоту \kod{z} сразу; загружается, но не летал & симулятор \\
\kod{circle(count, cell1, cell2, seconds, rotate='cw')} & круг вокруг середины клеток \kod{cell1} и \kod{cell2}; \kod{count} в полуоборотах; летал на дроне, время иное & проверено \\
\kod{circle_z(cell1, cell2, dz, seconds, rotate='cw')} & полкруга с подъёмом на \kod{dz} см; в дрон не загружается & симулятор \\
\kod{land()} & посадка на месте; в симуляции 3\,с & проверено \\
\kod{land_on(cell, seconds=2, low=35)} & посадка в центр клетки: подлёт, зависание 1,5\,с, снижение до \kod{low} за 2,5\,с, посадка & проверено \\
\bottomrule
\end{longtable}
\end{center}

\section*{Шоу}

\begin{center}
\small
\begin{tabular}{@{}p{4.9cm}p{9cm}@{}}
\toprule
Запись & Что делает \\
\midrule
\kod{Show(name)} & создаёт шоу; переменная должна называться \kod{show} \\
\kod{show.drone(number)} & возвращает дрона с этим номером, при первом обращении создаёт его \\
\bottomrule
\end{tabular}
\end{center}

\section*{Числа и пределы}

\begin{center}
\begin{tabular}{@{}lr@{}}
\toprule
Что & Значение \\
\midrule
клетка коврика & 40\cm \\
коврик & 12 $\times$ 12, клетки 0–143 \\
четверть & 6 $\times$ 6, 2,4 $\times$ 2,4\,м \\
лимит высоты по умолчанию & 150\cm (от 20 до 200\cm) \\
наименьшая дистанция между дронами & 60\cm \\
наибольшая средняя скорость перехода & 100\cm/с \\
посадка в симуляции & 3\,с \\
наибольшее ожидание на месте & 60\,с \\
батарея: предупреждение / запрет взлёта & 7,4\,В / 7,0\,В \\
наибольшее число шагов дрона в миссии & 255 \\
наибольшее число шагов в шоу & 10\,000 \\
время запуска кода шоу & 5\,с \\
\bottomrule
\end{tabular}
\end{center}
